%% Преамбула TeX-файла

% 1. Стиль и язык
\documentclass{G7-32} % Стиль (по умолчанию размер шрифта 14pt)

% Остальные стандартные настройки убраны в preamble.inc.tex.
\sloppy

% Настройки стиля ГОСТ 7-32
% Для начала определяем, хотим мы или нет, чтобы рисунки и таблицы нумеровались в пределах раздела, или нам нужна сквозная нумерация.
\EqInChapter % формулы будут нумероваться в пределах раздела
\TableInChapter % таблицы будут нумероваться в пределах раздела
\PicInChapter % рисунки будут нумероваться в пределах раздела

% Добавляем гипертекстовое оглавление в PDF
\usepackage[
bookmarks=true, colorlinks=true, unicode=true,
urlcolor=black,linkcolor=black, anchorcolor=black,
citecolor=black, menucolor=black, filecolor=black,
]{hyperref}

\AfterHyperrefFix

\usepackage{microtype}% полезный пакет для микротипографии, увы под xelatex мало чего умеет, но под pdflatex хорошо улучшает читаемость

% Тире могут быть невидимы в Adobe Reader
\ifInvisibleDashes
\MakeDashesBold
\fi

\usepackage{graphicx}   % Пакет для включения рисунков

% С такими оно полями оно работает по-умолчанию:
% \RequirePackage[left=20mm,right=10mm,top=20mm,bottom=20mm,headsep=0pt,includefoot]{geometry}
% Если вас тошнит от поля в 10мм --- увеличивайте до 20-ти, ну и про переплёт не забывайте:
\geometry{right=20mm}
\geometry{left=30mm}
\geometry{bottom=20mm}
\geometry{ignorefoot}% считать от нижней границы текста


% Пакет Tikz
\usepackage{tikz}
\usetikzlibrary{arrows,positioning,shadows}

% Произвольная нумерация списков.
\usepackage{enumerate}

% ячейки в несколько строчек
\usepackage{multirow}

% itemize внутри tabular
\usepackage{paralist,array}

%\setlength{\parskip}{1ex plus0.5ex minus0.5ex} % разрыв между абзацами
\setlength{\parskip}{1ex} % разрыв между абзацами
\usepackage{blindtext}

% Центрирование подписей к плавающим окружениям
\usepackage[justification=centering]{caption}

\usepackage{newfloat}
\DeclareFloatingEnvironment[
placement={!ht},
name=Equation
]{eqndescNoIndent}
\edef\fixEqndesc{\noexpand\setlength{\noexpand\parindent}{\the\parindent}\noexpand\setlength{\noexpand\parskip}{\the\parskip}}
\newenvironment{eqndesc}[1][!ht]{%
    \begin{eqndescNoIndent}[#1]%
\fixEqndesc%
}
{\end{eqndescNoIndent}}




% Настройки листингов.
\ifPDFTeX
% 8 Листинги

\usepackage{listings}

% Значения по умолчанию
\lstset{
  basicstyle= \footnotesize,
  breakatwhitespace=true,% разрыв строк только на whitespacce
  breaklines=true,       % переносить длинные строки
%   captionpos=b,          % подписи снизу -- вроде не надо
  inputencoding=koi8-r,
  numbers=left,          % нумерация слева
  numberstyle=\footnotesize,
  showspaces=false,      % показывать пробелы подчеркиваниями -- идиотизм 70-х годов
  showstringspaces=false,
  showtabs=false,        % и табы тоже
  stepnumber=1,
  tabsize=4,              % кому нужны табы по 8 символов?
  frame=single
}

% Стиль для псевдокода: строчки обычно короткие, поэтому размер шрифта побольше
\lstdefinestyle{pseudocode}{
  basicstyle=\small,
  keywordstyle=\color{black}\bfseries\underbar,
  language=Pseudocode,
  numberstyle=\footnotesize,
  commentstyle=\footnotesize\it
}

% Стиль для обычного кода: маленький шрифт
\lstdefinestyle{realcode}{
  basicstyle=\scriptsize,
  numberstyle=\footnotesize
}

% Стиль для коротких кусков обычного кода: средний шрифт
\lstdefinestyle{simplecode}{
  basicstyle=\footnotesize,
  numberstyle=\footnotesize
}

% Стиль для BNF
\lstdefinestyle{grammar}{
  basicstyle=\footnotesize,
  numberstyle=\footnotesize,
  stringstyle=\bfseries\ttfamily,
  language=BNF
}

% Определим свой язык для написания псевдокодов на основе Python
\lstdefinelanguage[]{Pseudocode}[]{Python}{
  morekeywords={each,empty,wait,do},% ключевые слова добавлять сюда
  morecomment=[s]{\{}{\}},% комменты {а-ля Pascal} смотрятся нагляднее
  literate=% а сюда добавлять операторы, которые хотите отображать как мат. символы
    {->}{\ensuremath{$\rightarrow$}~}2%
    {<-}{\ensuremath{$\leftarrow$}~}2%
    {:=}{\ensuremath{$\leftarrow$}~}2%
    {<--}{\ensuremath{$\Longleftarrow$}~}2%
}[keywords,comments]

% Свой язык для задания грамматик в BNF
\lstdefinelanguage[]{BNF}[]{}{
  morekeywords={},
  morecomment=[s]{@}{@},
  morestring=[b]",%
  literate=%
    {->}{\ensuremath{$\rightarrow$}~}2%
    {*}{\ensuremath{$^*$}~}2%
    {+}{\ensuremath{$^+$}~}2%
    {|}{\ensuremath{$|$}~}2%
}[keywords,comments,strings]

% Подписи к листингам на русском языке.
\renewcommand\lstlistingname{Листинг}
\renewcommand\lstlistlistingname{Листинги}

\else
\usepackage{local-minted}
\fi

% Полезные макросы листингов.
% Любимые команды
\newcommand{\Code}[1]{\textbf{#1}}


% Стиль титульного листа и заголовки
\usepackage{title-petrenko}

% --- Заполняем данные для титульника ---
% Те поля, которые вы не укажете, возьмут значения по умолчанию (БД, Петренко, ПНИПУ и т.д.)
\LabNumber{3} % Номер лабораторной
\LabTopic{Макет для отчетов работ} % Тема работы
\LabStudent{РИС-25-2б}{Ильинов Ф. М.} % Группа и ФИО студента

\begin{document}

\begin{titlepage}
\end{titlepage}

\frontmatter % выключает нумерацию ВСЕГО; здесь начинаются ненумерованные главы: реферат, введение, глоссарий, сокращения и прочее.

%\listoffigures                         % Список рисунков

%\listoftables                          % Список таблиц

%\NormRefs % Нормативные ссылки 
% Команды \breakingbeforechapters и \nonbreakingbeforechapters
% управляют разрывом страницы перед главами.
% По-умолчанию страница разрывается.

% \nobreakingbeforechapters
% \breakingbeforechapters


\tableofcontents

\printnomenclature % Автоматический список сокращений

\mainmatter % это включает нумерацию глав и секций в документе ниже


\chapter{Цель работы}

Реализовать вывод трёх SQL запросов (из лабораторной работы №2) на страницу браузера:
\begin{enumerate}
    \item Получить список сотрудников, номера телефонов и их заработные платы (листинг \ref{lst:sql1});
    \item Получить список сотрудников с их адресами (отсортировать по адресу) (листинг \ref{lst:sql2});
    \item Получить список сотрудников с продолжительностью трудовой деятельности больше 4 лет (листинг \ref{lst:sql3}).
\end{enumerate}

\begin{listing}[H]
    \sqlfile{../../lab02/query1.sql}
    \caption{SQL запрос для получения списка сотрудников, их номеров телефонов и их заработные платы}
    \label{lst:sql1}
\end{listing}

\begin{listing}[H]
    \sqlfile{../../lab02/query2.sql}
    \caption{SQL запрос для получения списка сотрудников с их адресами}
    \label{lst:sql2}
\end{listing}

\begin{listing}[H]
    \sqlfile{../../lab02/query3.sql}
    \caption{SQL запрос для получения списка сотрудников, чья трудовая деятельность больше 4 лет}
    \label{lst:sql3}
\end{listing}
\chapter{Задачи работы}
\label{cha:tasks}

В соответствии с целью были поставлены следующие задачи:

\begin{enumerate}
    \item Спроектировать и разработать локальный веб-сервер;
    \item Спроектировать и разработать веб-интерфейс;
    \item Вывести результаты запросов на страницу браузера.
\end{enumerate}
\chapter{Этапы выполнения}

Для выполнения поставленных задач были сформулированы следующие технологические требования к инструментарию:
\begin{itemize}
    \item Поддержка работы с СУБД PostgreSQL \cite{doc:postgres};
    \item Поддержка работы с протоколом HTTP \cite{wiki:http};
    \item Возможность отображения кнопок в веб-интерфейсе;
    \item Возможность отображения таблиц в веб-интерфейсе.
\end{itemize}

Для работы с веб-страницами был выбран язык разметки HTML \cite{wiki:html}, так как он позволяет создавать кнопки, требуемые
для выбора текущего SQL запроса, и отрисовывать таблицы.

В качестве языка программирования веб-сервера был выбран Golang \cite{wiki:go}, так как его стандартные библиотеки реализовывают широкий
и удобный функционал для работы с веб-серверами.

Программа реализует принцип Server-Side Rendering.
Итоговые HTML собираются на стороне сервера и отправляются клиенту.

\backmatter %% Здесь заканчивается нумерованная часть документа и начинаются ссылки и

\Conclusion % заключение к отчёту

В процессе выполнения работы были освоены навыки работы с языком разметки HTML,
работа с СУБД PostgreSQL в языке Golang, работа с протоколом HTTP в языке Golang.

Также был реализован вывод результатов трех SQL запросов на страницу браузера:
\begin{enumerate}
    \item Получить список сотрудников, номера телефонов и ЗП (листинг \ref{lst:sql1});
    \item Получить список сотрудников с их адресами (отсортировать по адресу) (листинг \ref{lst:sql2});
    \item Получить список сотрудников с продолжительностью трудовой деятельности больше 4 лет (листинг \ref{lst:sql3}).
\end{enumerate}

%%% Local Variables: 
%%% mode: latex
%%% TeX-master: "rpz"
%%% End: 

% % Список литературы при помощи BibTeX
% Юзать так:
%
% pdflatex rpz
% bibtex rpz
% pdflatex rpz

\bibliographystyle{ugost2008}
\bibliography{report-lab-03}

%%% Local Variables: 
%%% mode: latex
%%% TeX-master: "rpz"
%%% End: 


\appendix   % Тут идут приложения

\end{document}

%%% Local Variables:
%%% mode: latex
%%% TeX-master: t
%%% End:
